\subsection{Decay of the archimedean projector}\label{sec:count-decay}

The finite projector restricts rational elements. At the real places the spectral gaps must give uniform decay away from the compact subgroup.

At one real place use the coordinate order of the strictly decreasing
Harish--Chandra parameter \(\lambda=(\lambda_1,\ldots,\lambda_r)\).
Let \(J\) be the diagonal signature matrix in this order, and
\(K=\mathrm U(p)\times\mathrm U(q)\) its compact centralizer in
\(\mathrm U(r)\). The gap condition is
\(\lambda_i-\lambda_j\ge h\) for \(i<j\) of opposite signs.
This allows either sign at each index in the weight order.

Write \(S\) for the irreducible lowest \(K\)-type, \(w\) for a unit
highest weight vector of \(S\), and \(P_z\) for the rank-one
orthogonal projector onto a unit vector \(z\). For a Hermitian
matrix \(X\) with \(XJ=-JX\), the compression of \(\pi(\exp X)\) to
\(S\) is
\begin{equation}\label{eq:count-compression}
 (\dim S)\int_K \Phi_\lambda(k^{-1}\exp(X)k)P_{\pi(k)w}\,dk .
\end{equation}
The Haar measure on \(K\) has mass one. Let \(\rho=\rho_C+\rho_N\) be the half-sum of the positive roots,
split into compact and noncompact roots.
If \(m_i(A)\) denotes the
leading principal minor of order \(i\) and \(m_0=1\), define
\[
 \Phi_\lambda(A)=\prod_{i=1}^r \left(\frac{m_i(A)}{m_{i-1}(A)}\right)^{-\mu_i},
 \qquad \mu=\lambda+\rho.
\]
The exponents \(\mu_i\) are
integers by the Harish--Chandra parameter condition. For \(A\) in
\eqref{eq:count-compression}, \(m_r(A)=1\) since
\(\operatorname{tr} X=0\). Thus only the powers
\(-\mu_i+\mu_{i+1}=-(\lambda_i-\lambda_{i+1}+1)\) for \(i<r\) remain.

\begin{proposition}[Lowest-type compression and polar decay]\label{prop:lowest-type-polar-decay}
Let \(\pi\) be the discrete series of
Definition~\ref{hyp:cuspidal-dimension}, with lowest \(K\)-type \(S\),
and let \(X=X^*\) satisfy \(XJ=-JX\). For \(q\ge1\), the lowest-type compression is the flag integral
\eqref{eq:count-compression}. Writing \(P_S\) for the lowest-type projector, its norm bound \eqref{eq:count-polar-decay} is 
\[
 \|P_S\pi(\exp X)|_S\|_{\rm op}\le\exp\!\left[-\frac{h}{8}\min\!\left(s^2,\frac{s}{\sqrt{2q}}\right)\right],\qquad s=\|X\|_{\rm HS}.
\]

If \(q=0\), the relation \(XJ=-JX\) forces \(X=0\) and the
compression is the identity.
\end{proposition}

\begin{proof}
\proofstep{The compact flag integral.}
Let \(H(A)\) denote the expression on the right in
\eqref{eq:count-compression}, with \(A\) in place of \(\exp X\).
It is defined on positive definite matrices and is holomorphic
where all the minors in the integral are uniformly nonzero,
in particular near the identity. Write \(\tau\) for the
holomorphic extension of the action on \(S\) to the complexified
compact group \(K_{\C}\). It satisfies \(\tau(a^*)=\tau(a)^*\),
where the matrix star in this paragraph denotes Hermitian transpose.
This holds on \(K\) by unitarity and extends to \(K_{\C}\) by
holomorphic continuation. The highest weight and the minor exponent satisfy
\[
 \Lambda=\lambda+\rho_N-\rho_C,\qquad
 \mu=\lambda+\rho_N+\rho_C=\Lambda+2\rho_C .
\]
This difference of weights determines the Jacobian in the flag integral.

For \(a\in K_{\C}\) and \(k\in K\), apply QR decomposition in each
sign space using the induced coordinate order. Write \(ak=cb\),
where \(c\in K\) and \(b\) is upper triangular with positive diagonal
\(d_i\). Write \(b_i\) for the first \(i\)-by-\(i\) block of \(b\), and
\(A_i\) for that of \(A\). Triangularity implies
\[
 (b^*Ab)_i=b_i^* A_i b_i,\qquad
 m_i(b^*Ab)=|\det b_i|^2m_i(A)=\prod_{j\le i}d_j^2\,m_i(A).
\] Since the upper unipotent fixes the highest weight vector, the
minor and vector factors transform by
\[
 \begin{aligned}
 \Phi_\lambda(k^{-1}a^*Aa k)&=d^{-2\mu}\Phi_\lambda(c^{-1}Ac),\\
 \tau(k)w&=d^\Lambda\tau(a)^{-1}\tau(c)w,
 \end{aligned}
 \qquad d^\nu=\prod_i d_i^{\nu_i}.
\] The map of compact flags \(kT\mapsto cT\), for the diagonal
torus \(T\) of \(K\), has real Jacobian \(d^{-4\rho_C}\).

In the tangent spaces translated to the standard flag, its complex
differential is \(\operatorname{Ad}(b)\) on the quotient of the
compact complexified Lie algebra by its upper triangular subalgebra.
The quotient has one direction for each negative compact root, and
the diagonal of the induced triangular action is \(d_j/d_i\) for
\(i<j\) of the same sign. Squaring the absolute value of the
determinant gives the stated Jacobian in the invariant flag measure.
The integral descends to this flag space because a torus element
preserves the highest weight line and conjugation by it preserves
all the leading minors.

Inserting the two weights in the integral
gives precisely \(d^{-2\mu+2\Lambda}=d^{-4\rho_C}\). Therefore
\[
       H(a^* A a)=\tau(a)^{-1}H(A)\tau(a^*)^{-1} .
\]
Near identity this identity extends holomorphically in \(A\) and in
the independent variables replacing \(a^*\) and \(a\). Indeed the
Hermitian matrices and the pairs \((a^*,a)\) give real forms in
local complex coordinates for those variables. The local holomorphic identity is therefore
\[
 H(uAv)=\tau(v)^{-1}H(A)\tau(u)^{-1}
       \qquad(u,v\in K_{\mathbb C}\text{ near }I).
\]
It follows that the germ \(g\mapsto H(g^*)\) on the real unitary
group has the left and right \(K\)-covariance of the compression.
Its value at identity is \(I_S\) by compact Schur orthogonality.

\proofstep{The first-order equations.}
The complexified noncompact tangent space \(\mathfrak p_{\C}\) has
basis \(E_{ij}\) with opposite signs. Its tensor product with \(S\)
splits into two orthogonal sums \(V_+\) and \(V_-\) of compact types
with highest weights \(\Lambda+\alpha\) for positive and negative
noncompact roots \(\alpha\), respectively. To see the list of types,
let \(W_+,W_-\) be the two sign spaces. Under the complexified
compact group,
\[
 \mathfrak p_{\C}=(W_+\otimes W_-^*)\oplus(W_-\otimes W_+^*).
\]
The matrix unit from the \(j\)-th to the \(i\)-th coordinate
has weight \(e_i-e_j\), with the signs of \(i,j\) opposite. The one-box Pieri rule and its dual add a single
coordinate in the first factor and subtract a single coordinate in
the other, retaining exactly the dominant results. These are the
indicated noncompact roots.

In the discrete series every compact type
has highest weight \(\Lambda\) plus a sum of positive roots for the
order of \(\lambda\), by Atiyah--Schmid, Theorem 8.5. Thus the action
map \(\mathfrak p_{\C}\otimes S\to\mathcal H_\pi^\infty\) kills
\(V_-\). This gives the first-order equations for the compression
\(g\mapsto P_S\pi(g)|_S\). For \(v_-=\sum E_{ij}\otimes u_{ij}\in V_-\), differentiate
the compression \(F(g)=P_S\pi(g)|_S\) in the right directions:
\[
 \sum D_{ij}F(g)u_{ij}
   =P_S\pi(g)\left(\sum d\pi(E_{ij})u_{ij}\right)=0 .
\]
The vanishing is precisely the absence of \(V_-\) in the
discrete-series action on \(\mathfrak p_{\C}\otimes S\).

For an upper strictly triangular \(Y\), right multiplication
by \(\exp(tY)\) acts on each initial block of columns by an
upper unitriangular matrix. In particular,
\[
 m_l(Z\exp(tY))=m_l(Z)\quad(1\le l\le r),\qquad D_Y\Phi_\lambda=0 .
\] The scalar in \(H(g^*)\) is
\(\overline{\Phi_\lambda(k^{-1}gk)}\). On the real Lie algebra
\(\overline Y=-JY^tJ\), so after complexification and conjugation
by \(k\) its derivative only uses the original entries \(Y_{ij}\)
with \(i\le j\). The conjugate of each vanishing upper-direction
derivative is a vanishing lower-direction derivative.
In the tensor of differentiation directions and input vectors this
conjugation is the compact action by \(k^{-1}\), which preserves
\(V_-\). For \(i<j\), the component indexed by the matrix unit \(E_{ij}\) has zero coefficient along \(w\). Otherwise that tensor would have a weight component
\(\Lambda+e_i-e_j\).

Indeed every weight of a constituent with highest weight
\(\Lambda-\beta\), for a positive noncompact root \(\beta\), has the form
\[
 \Lambda-\beta-\sum_{\alpha\in\Delta_C^+} n_\alpha\alpha,
       \qquad n_\alpha\in\mathbb N ,
\]
where \(\Delta_C^+\) is the set of positive compact roots. This
omits \(\Lambda+e_i-e_j\) for \(i<j\), since a nonempty sum of
positive roots cannot vanish. 

The diagonal matrix units do not belong to \(\mathfrak p_{\C}\). After conjugation by \(k\), write the input tensor as \(\sum_{i,j} E_{ij}\otimes u_{ij}(k)\), summing over the noncompact root spaces. The differentiated integrand
is a sum of scalar derivatives times
\(\langle w,u_{ij}(k)\rangle\pi(k)w\).
If \(a_{ij}(k,g)\) denotes the scalar derivative in the direction \(E_{ij}\), the
two tests give
\[
 a_{ij}(k,g)=0\quad(i>j),\qquad
 \langle w,u_{ij}(k)\rangle=0\quad(i<j).
\]
There is no diagonal term, so each summand
\(a_{ij}(k,g)\langle w,u_{ij}(k)\rangle\pi(k)w\) vanishes. Thus the first-order equation holds before integration.

\proofstep{Determination from the value at identity.}
Let \(F\) be the difference of the two germs. It is zero at identity.
Suppose derivatives of all orders less than \(m\) vanish there.
At order \(m\) any compact differentiation can be commuted to the
right, with commutators of lower order. Right covariance then
reduces it to order \(m-1\). Interchanging two noncompact directions produces a commutator of
lower order, which vanishes by induction. Thus the remaining derivative belongs to
\[
 \operatorname{Hom}_K(\operatorname{Sym}^m\mathfrak p_{\mathbb C}
                         \otimes S,S).
\] Differentiating the
first-order equations shows the map kills tensors with \(V_-\)
in the product of any one noncompact tensor factor with \(S\). It factors through
the orthogonal complement of the sum of these subspaces. In this
complement each noncompact factor paired with \(S\) lies in \(V_+\).

Let \(c_S\) be the compact Casimir eigenvalue on \(S\), in the trace
form. On the tensor product the compact Casimir is
\(c_S+\Omega_{\mathfrak p^{\otimes m}}+2\sum_i M_i\), where
\(M_i\) is the mixed term between the \(i\)-th tensor factor and \(S\).
Compact unitarity gives
\[
 \langle u,\Omega_{\mathfrak p^{\otimes m}}u\rangle\ge0 .
\] The Casimir on a single \(\mathfrak p_{\C}\)
factor equals \(p+q=r\), since the standard and dual standard
representations of \(\mathrm U(d)\) each have trace-form Casimir
\(d\). Thus on a positive type \(\Lambda+e_a-e_b\), \(2M_i\)
has eigenvalue
\[
  2(\lambda_a-\lambda_b+\rho_N(a)-\rho_N(b))+2-r
                 \ \ge\ 2(h-2r).
\]
Here \(|\rho_N(a)|\le r/2\).

Since \(h>2r\) in the noncompact
case, the quadratic form of the tensor Casimir on the complement
is strictly larger than \(c_S\) times the squared norm. The
complement is \(K\)-invariant, so it has no summand isomorphic to
\(S\). This forces the derivative map to vanish and completes the
induction.

Along \(g=\exp(tX)\) with \(X\) Hermitian and anticommuting with
\(J\), the integral is real analytic for all real \(t\). All the
leading minors are positive there and compactness gives a common
complex neighborhood of each \(t\) for the integration. The
compression is real analytic on this curve as well. In fact the
infinitesimal norm estimate \eqref{eq:count-step} below at one
place grows at most linearly in the word length. Iteration from the lowest type gives a constant \(C_X\) such that
\[
 \left\|\frac{d^m}{dt^m}\pi(\exp(tX))u\right\|
       \le m! C_X^m\|u\|\qquad (m\ge0,\ t\in\mathbb R).
\]
Unitarity makes the bound uniform in \(t\), so Taylor's estimate
gives analyticity. The vanishing
at identity of all derivatives therefore proves
\eqref{eq:count-compression} for every such \(X\). The quantitative threshold \(h\ge120rq^{3/2}\) implies
\(h>2r\) at every noncompact place.

\proofstep{Quantitative decay.}

The local expression for the compressed coefficient is now explicit. Its useful decay must be expressed in terms of the opposite-sign spectral gap, which remains large throughout the count.


Put \(Y=k^{-1}Xk\) and \(s=\|Y\|_{\mathrm{HS}}=\|X\|_{\mathrm{HS}}\).
On the \(i\)-th exterior power let \(L_i(Y)\) act by applying
\(Y\) to each wedge factor and summing. With
\(e_{[i]}=e_1\wedge\cdots\wedge e_i\), the minor has the spectral representation
\[
 m_i(\exp Y)=\langle e_{[i]},\exp(L_i(Y))e_{[i]}\rangle.
\]

The spectral measure of \(L_i(Y)\) at this unit vector is a
probability distribution \(\xi_i\). Its mean is zero because all
diagonal entries of \(Y\) vanish. Its variance is
\[
   V_i=\|L_i(Y)e_{[i]}\|^2=\sum_{a\le i<b}|Y_{ba}|^2.
\]

Indeed, for each term off the diagonal, the replacement of the
\(a\)-th factor by \(e_b\) either repeats a factor or gives one of
the orthonormal wedges with exactly one index outside \([i]\).
The eigenvalues of \(Y\) are pairs \(\pm t_j\) with \(t_j\ge0\),
and zeros, by the singular-value decomposition of the block between
the two sign spaces. There are at most \(q\) pairs and
\(2\sum_j t_j^2=s^2\). Each eigenvalue of \(L_i(Y)\) is a sum drawn without replacement
from this multiset. In particular
\[
              \|L_i(Y)\|_{\rm op}\le \sum_j t_j\le\sqrt{2q}\,s .
\]

For \(s>0\), set \(t=\min(1,1/(\sqrt{2q}\,s))\). If \(|z|\le1\),
Taylor's integral remainder and \(e^{uz}\ge1-u\) for \(0\le u\le1\)
give \(e^z\ge1+z+z^2/3\). Apply this to \(z=t\xi_i\). Since
\(t^2V_i\le1\) and \(\log(1+a/3)\ge a/4\) for \(0\le a\le1\),
the log moment satisfies \(\log\mathbb E e^{t\xi_i}\ge t^2 V_i/4\).
Its value at zero is zero, so convexity on \([0,1]\) gives
\[
 \log m_i(\exp Y)=\log\mathbb E e^{\xi_i}
       \ge t^{-1}\log\mathbb E e^{t\xi_i}\ge tV_i/4 .
\] The weighted sum of
the variances satisfies
\[
 \sum_{i<r}(\lambda_i-\lambda_{i+1}+1)V_i
 = \sum_{a<b}(\lambda_a-\lambda_b+b-a)|Y_{ba}|^2
 \ge h s^2/2 .
\]
All nonzero entries in the last sum cross the two signs. Thus the
scalar multiplier in \eqref{eq:count-compression} is bounded above by
\begin{equation}\label{eq:count-polar-decay}
      \exp\!\left[-\frac h8\min\!\left(s^2,\frac{s}{\sqrt{2q}}\right)\right].
\end{equation}

It is nonnegative, and the same bound is immediate for \(s=0\).
By Schur orthogonality,
\((\dim S)\int_K P_{\pi(k)w}\,dk=I_S\). Consequently the multiplier
bound is also the operator norm bound for the compression, with no
factor depending on \(\dim S\). The polar decomposition \(g=k_0\exp X\)
preserves this bound since \(\pi(k_0)\) acts unitarily on \(S\).

\end{proof}

\subsection{The initial bound in the deep part}\label{sec:count-sobolev}

Near a very short lattice, many translates can contribute at a single point. A pointwise bound for any finite orthonormal family is needed before one knows that the whole projected space has finite dimension.

\begin{proposition}[Initial point evaluation in a fixed order]\label{prop:initial-point-evaluation}
Fix the archimedean representation, Haar gauges and finite level of
Definition~\ref{hyp:cuspidal-dimension}, with \(r\ge8\).
Set \(N=r^2 f\), \(B=b_{P_0+N}\) as in \eqref{eq:count-step}, and
let \(C(y)\) be the multiplicity factor in \eqref{eq:gram-chart-multiplicity}.
For a finite family of \(D\) isometric intertwiners and a vector \(u\)
of order at most \(P_0\) at each place, one has
\[
 \sum_{j=1}^D |F_j(u,y)|^2\le (4B)^N C(y)D\|u\|^2 .
\]
For \(P_0=8qN\) this gives the starting bound, uniform on the
neighborhood of a deep point with least size \(\tau(x)\), used
in the descent.
\end{proposition}

\begin{proof}
Use the same isometric intertwiners that give the diagonal kernel.
Write \(F_j(u,y)\) for the continuous automorphic representative at
\(y\) of a vector \(u\) in the \(j\)-th copy of the representation.
It is invariant on the right by the finite level subgroup \(K_g\),
and \(\|F_j(u,\cdot)\|_2=\|u\|\). At each real place take the
filtration by the span of words of at most \(P\) infinitesimal
operators on the lowest \(K\)-type. Use the tensor product of these
spaces, with the same bound \(P\) at each place, as the space \(V_P\).
For a unit Hilbert--Schmidt Lie algebra element at one place,
differentiation sends \(V_P\) into \(V_{P+1}\) with norm at most
\begin{equation}\label{eq:count-step}
 b_P=\sqrt r\,b^2+\tfrac32 r^{3/2}+3P .
\end{equation}

To see this, let \(\rho_C,\rho_N\) be the half-sums of positive
compact and noncompact roots, and \(\rho=\rho_C+\rho_N\). The
lowest \(K\)-type has highest weight \(\Lambda=\lambda+\rho_N-\rho_C\).
For the derived group this is the discrete-series lowest-type formula
of Atiyah and Schmid, Theorem 8.5. The central character supplies
the common central coordinate when passing from
\(\mathrm{SU}(p,q)\times \mathrm U(1)\) to its finite central quotient
\(\mathrm U(p,q)\).
The infinitesimal character is \(\lambda\) in the normalization in
which the quadratic Casimir for the trace form acts by
\(\|\lambda\|^2-\|\rho\|^2\). \cite{atiyah-schmid}
The lowest isotypic component consists of smooth vectors.

Write \(P_S\) for the compact isotypic projection. Compact
averaging preserves smoothness, and continuity gives
\[
 P_S(\mathcal H_\pi^\infty)\subseteq S\cap\mathcal H_\pi^\infty,
 \qquad \overline{P_S(\mathcal H_\pi^\infty)}=P_S(\mathcal H_\pi)=S .
\]
Since \(S\) is finite dimensional, the projected dense linear space
equals \(S\).
For the word space \(V_P\), the surjection
\[
 \bigoplus_{j=0}^P\mathfrak g_{\C}^{\otimes j}\otimes S
      \longrightarrow V_P,\qquad
 X_1\otimes\cdots\otimes X_j\otimes u
       \longmapsto d\pi(X_1)\cdots d\pi(X_j)u
\]
is \(K\)-equivariant.
Thus each occurring compact highest weight \(\nu\) is a weight of
the first factor plus at most \(P\) roots of length \(\sqrt2\).
The weights of the first factor lie in the convex hull of the compact
Weyl orbit of \(\Lambda\). For any compact Weyl permutation \(w\), dominance gives
\[
 \|\Lambda+\rho_C\|^2-\|w\Lambda+\rho_C\|^2
      =2\langle\Lambda-w\Lambda,\rho_C\rangle\ge0 .
\]
Convexity extends the norm bound to every weight in the convex hull. Adding the contributions from the lowest type and the \(P\) adjoint
factors yields
\[
         \|\nu+\rho_C\|\le\|\lambda\|+\|\rho_N\|+\sqrt2 P .
\]

Let \((E_i)\) be a Hilbert--Schmidt orthonormal basis
adapted to the compact and noncompact directions. In the trace
form the compact basis contributes with a minus sign and the
noncompact basis with a plus sign. Thus in the enveloping algebra
\[
 \Omega_K=-\sum_{i\ {\rm compact}}E_i^2,\qquad
 \Omega_G=\Omega_K+\sum_{i\ {\rm noncompact}}E_i^2.
\]
Since \(d\pi(E_i)\) is skew-adjoint on smooth vectors, it follows
that \(\sum_i\|d\pi(E_i)u\|^2
=\langle u,d\pi(2\Omega_K-\Omega_G)u\rangle\). The operator \(\Omega_K\) uses the trace form
restricted to the compact subalgebra and acts on a compact type
\(\nu\) by \(\|\nu+\rho_C\|^2-\|\rho_C\|^2\).

With
\(A=\|\lambda\|\) and \(\delta=\|\rho_N\|+\sqrt2 P\), the sum is
therefore at most \((2(A+\delta)^2-A^2+\|\rho\|^2)\|u\|^2\).
In fact
\[
 2(A+\delta)^2-A^2+\|\rho\|^2
        \le (A+2\delta+\|\rho\|)^2,\qquad
 A+2\delta+\|\rho\|
        \le \sqrt r\,b^2+\tfrac32r^{3/2}+2\sqrt2 P\le b_P .
\]
The second comparison uses the bound \(r/2\) on each coordinate
of \(\rho,\rho_N\) and \(b^2\) on \(|\lambda_i|\). For an arbitrary unit direction use
Cauchy--Schwarz in the orthonormal basis, and tensor the resulting
bound at a single place with the identity at the others.

\proofstep{An injective product chart and its Jacobian.}

Write \(N=r^2 f\) and \(B=b_{P_0+N}\) for a fixed starting order
\(P_0\). Take \(a=1/B\). By Definition~\ref{hyp:cuspidal-dimension}, \(b^2\ge r^{19}\);
for \(r\ge8\) the chart radius therefore satisfies
\[
   a=B^{-1}\le\min\{(8r^5)^{-1},(12r^3)^{-1}\}.
\]
At each real place use an orthonormal Lie algebra basis \(E_i\)
and the coordinates
\(t\mapsto\Theta(t)=\prod_{i=1}^{r^2}\exp(t_i E_i)\) on
\([0,a]^{r^2}\). For a product \(U=\prod_j\exp(s_j E_{i_j})\), put \(L=\sum_j|s_j|\); negative parameters include inverses. Expanding the product and using \(\|E_i\|_{\rm op}\le1\) gives 
\[
 \|U\|_{\rm op}\le\prod_j e^{|s_j|}=e^L,\qquad \|U-I\|_{\rm op}\le\prod_j e^{|s_j|}-1=e^L-1.
\]

For the chart differential, its \(i\)-th column is a product of a
prefix, \(E_i\), and a suffix. The ambient Hilbert--Schmidt distance
from \(E_i\) is at most \(e^{r^2 a}-1\le 2r^2a\). Let \(S_i=\prod_{j>i}\exp(t_jE_j)\). Left translation gives
\[
 \Theta(t)^{-1}\partial_i\Theta(t)=S_i^{-1}E_i S_i,\qquad
 \|S_i^{-1}E_iS_i-E_i\|_{\rm HS}\le e^{2r^2a}-1\le3r^2a .
\]
Indeed subtract \(E_iS_i\) and use the bounds on \(S_i^{\pm1}\)
together with \(\|E_i\|_{\rm HS}=1\).
The operator norm of a map with \(r^2\) such error columns is at
most \(r\) times the maximum column norm. Set \(I_E(t)=\sum_i t_iE_i\), the linear map into the
ambient matrix space. For \(t,t'\) in the cube, integration of
the ambient differential along the segment gives
\[
 \|\Theta(t)-\Theta(t')-I_E(t-t')\|_{\rm HS}
                         \le 2r^3a\|t-t'\| .
\]

Since \(\|I_E(t-t')\|_{\rm HS}=\|t-t'\|\)
and \(2r^3a<1\), the chart is injective. For the
left-translated differential each singular value is at least
\(1-3r^3a\), so the Jacobian with respect to the orthonormal basis satisfies
\[
 |\det d\Theta|\ge(1-3r^3a)^{r^2}\ge e^{-6r^5a}\ge e^{-1}.
\]
Here \(3r^3a\le1/4\) and \(6r^5a\le1\).
Let \(c_v\) be the reciprocal of the density of Haar volume at the
identity relative to Hilbert--Schmidt Lebesgue measure. Write \(dt_v\) for coordinate Lebesgue measure and \(dg_v\) for
Haar measure. On the image the Jacobian estimate gives
\[
 dt_v\le e c_v\,dg_v,\qquad
 \prod_v dt_v\le e^f\Big(\prod_v c_v\Big)\prod_v dg_v .
\] Two elements of the image
have quotient of operator norm \(\le e^{2r^2 a}\le 2\) at every place.

\proofstep{Unfolding the \(L^2\) integral.}

Write \(\tau_1(y)=\min(1/2,\tau(y))\) for the truncated least
lattice size at \(y\). In unfolding the chart about \(y\) with the
finite coset \(y_fK_g\), two points mapping to the same automorphic
point differ by a rational group element preserving the lattice at
\(y\). Relative to a chosen representative of the fibre, its
archimedean matrix is a quotient of two chart points. Thus
\[
 \|\gamma_v\|_{\rm HS}\le\sqrt r\|\gamma_v\|_{\rm op}
     \le\sqrt r e^{2r^2a}\le2\sqrt r
\]
in the lattice frame at each place.
The lattice separation and product packing in the far orbital count
below bound the number in the fibre by
\(M(y)=(1+4\sqrt r/\tau_1(y)^2)^{2N}\).
Indeed the difference of two distinct elements is a nonzero integral
endomorphism, so the geometric mean of its Hilbert--Schmidt norms
is at least \(\tau_1(y)^2\).

Product balls about the elements with
radius \(b_0=\tau_1(y)^2/2\) at each place thus have disjoint interiors.
In ambient real dimension \(2N\) their volume ratio is bounded by
\[
 \left(\frac{2\sqrt r+b_0}{b_0}\right)^{2N}
     =\left(1+\frac{4\sqrt r}{\tau_1(y)^2}\right)^{2N}.
\] The finite coset contributes its full
volume \(\operatorname{vol}(K_g)\) in the unfolded integral, since
the squared absolute value is right \(K_g\)-invariant. The integral
of \(|F_j(u,\cdot)|^2\) on the product coordinate cube is therefore
at most \(C(y)\|u\|^2\), where
\begin{equation}\label{eq:gram-chart-multiplicity}
 C(y)=e^f\Bigl(\prod_v c_v\Bigr)\frac{M(y)}{\operatorname{vol}(K_g)} .
\end{equation}
This estimate uses the isometry and applies to every
\(u\in V_{P_0+N}\).

\proofstep{Iteration of the point evaluation inequality.}

The evaluation bound has so far been obtained at a fixed derivative order. Iteration allows the depth of the thin region to determine how much decay is needed.


For a differentiable function \(\phi\) on \([0,a]\), the fundamental
theorem of calculus and Cauchy--Schwarz give
\[
 |\phi(0)|^2\le\frac{2}{a}\int_0^a|\phi(t)|^2\,dt
                     +2a\int_0^a|\phi'(t)|^2\,dt .
\]
Start at the origin and append the coordinates in order. At stage
\(i\) fix the prefix \(x=y\exp(t_1 E_1)\cdots
\exp(t_{i-1}E_{i-1})\) and one of the fixed vectors \(w\) from
the preceding stages.

Apply the one-variable inequality to
\(t\mapsto F_j(w,x\exp(tE_i))\). Right differentiation inserts the next word operator:
\[
 \frac d{dt}F_j(w,x\exp(tE_i))=F_j(d\pi(E_i)w,x\exp(tE_i)).
\] The two resulting
cube integrands thus use the fixed vectors \(w\) and
\(d\pi(E_i)w\), with weights \(2/a\) and \(2a\),
respectively. Applying the next inequality at each point of this
prefix leaves these vectors fixed as the next factor is appended.
After \(N\) stages apply the cube bound to each resulting vector.
The sum of their squared norms with these weights is at most
\((2/a+2aB^2)^N\|u\|^2\), so for \(u\in V_{P_0}\)
\[
 \sum_{j=1}^D |F_j(u,y)|^2\le(2/a+2a B^2)^N C(y)D\|u\|^2
                          =(4B)^N C(y)D\|u\|^2 .
\]
Each derivative raises the order by at most one, so \(P_0+N\)
suffices for the whole iteration. For the subsequent descent one may enlarge \((4B)^N\) to \((6B)^N\).

In the unipotent neighborhood used below the finite lattice is
fixed. The archimedean change of frame at each place is
\(\exp(A_v)\exp(Z_v)\), with
\(\|A_v\|_{\mathrm{HS}}+\|Z_v\|_{\mathrm{HS}}\le\rho\le\log 2\).
The matrix and its inverse both have operator norm at most two.
Consequently each covolume in rank \(d\) loses at most a factor
\(2^d\) at one place in the complex Euclidean determinant norm.
For a fixed subquotient of rank \(d>0\) the inverse frame has
\(\|\bigwedge^d g_v^{-1}\|\le2^d\). The induced quotient norm obeys
the same bound by taking the infimum over lifts. Thus
\[
  C_y/C_x\ge\prod_{v\mid\infty}2^{-2d}=2^{-2fd},\qquad
  (C_y/C_x)^{1/(2fd)}\ge\tfrac12
\]
for its normalized covolumes.

Taking the least and truncating at \(1/2\) shows
\(\tau_1(y)\ge\tau(x)/2\) when \(\tau(x)<1/2\). There
\(1+4\sqrt r/\tau_1(y)^2\le17\sqrt r/\tau(x)^2\). The factor for
that neighborhood follows with \(P_0=8qN\). These bounds use the
given lattice and gauge and hold uniformly in \(r,f,b,P_0\) under
the stated size conditions.
\end{proof}
